%قبل از شروع مسئله ترجیحاً چیزی ننویسید. اما اگر خواستید بنویسید، در قالب لیستی (فقرات و شمارش) بنویسید.
%برای اضافه کردن نمره، بعدی از شروع‌مسئله، [x] را اضافه کنید و به جای x  نمره مورد نظر را قرار دهید. برای زیرمسئله‌ها هم به همین صورت است.
%\شروع‌مسئله[10]
\شروع‌مسئله
% صورت اصلی مسئله را در این‌جا بنویسید.
میخواهیم روی یک رشته $n$ بیتی از ۰ و ۱ الگوریتم زیر را اجرا کنیم. با توجه به این الگوریتم به سوالات زیر پاسخ دهید.

\lr{
\begin{algorithm}
	\begin{algorithmic}
		\While{$true$} 
		\State $cnt := 0$
		\For{$ i = 1 \to n-3$}
		\If{$s[i] = 1 \; and \; s[i+1] = 0 \; and \; s[i+2] = 1 $}
		\State $s[i] := 1 - s[i]$
		\State $s[i+1] := 1 - s[i + 1]$
		\State $s[i+2] := 1 - s[i+2]$
		\State $cnt := cnt + 1$
		\EndIf
		\EndFor
		\If{$cnt = 0$}
		\State $break$
		\EndIf
		\EndWhile  \label{roy's loop}
	\end{algorithmic}
\end{algorithm}
}






%اگر مسئله چند قسمت داشت به این صورت عمل کنید:
\شروع{زیرمسائل}
\زیرمسئله
ثابت کنید این الگوریتم به پایان می‌رسد.
\زیرمسئله
پیچیدگی زمانی الگوریتم را در بهترین و بدترین حالت حساب کرده و برای هر کدام یک مثال بزنید.
%\زیرمسئله % مسئلۀ ۱ت
\پایان{زیرمسائل}


\solution{
برای اتمام پایان پذیری این الگوریتم، می‌خواهیم از نزول متناهی استفاده کنیم. تعداد یک‌های رشته‌ی اولیه را در نظر بگیرید. هر بار که if اجرا شود، تعداد یک‌ها یکی کاهش می‌یابد. در نتیجه تعداد دفعاتی که if اجرا می‌شود محدود است. پس بالاخره در یک زمان، for اجرا می‌شود بدون این کهif اجرا شود و الگوریتم در آن مرحله اتمام می‌پذیرد.

می دانیم در هر صورت حداقل یک بار for اجرا می‌شود. پس حداقل اردر زمانی
$O(n)$
است. هم‌چنین اگر همه‌ی رشته از ۰ تشکیل شده باشد، تنها یک بار for اجرا می‌شود و در نتیجه الگوریتم از اردر 
$\Theta(n)$
خواهد بود.

اما ادعا می‌کنیم در بدترین حالت، این الگوریتم از اردر 
$\Theta(n^2)$
خواهد بود. برای این حالت، رشته‌ی زیر را در نظر بگیرید که از 
$n-1$
تا ۱ و ۱ عدد ۰ تشکیل شده است.
$$1111 ... 1101$$
پس از اجرای هر بار for، تنها زیر رشته‌ی 101 یکی به سمت چپ شیفت می‌خورد. بنابراین در اردر 
$\Theta(n)$
بار for اجرا می‌شود . هر بار اجرا هم از اردر 
$\Theta(n)$
طول می‌کشد. پس اردر کل الگوریتم از اردر 
$\Theta(n^2)$
خواهد شد.

}




%%%%%%%%%%%%%%%%%%%%%%%%%%%%%%%%%%%%%%%%%%%%%%%%%%%%%
% برخی دستورات مفید برای تمرین‌ها. %
%%%%%%%%%%%%%%%%%%%%%%%%%%%%%%%%%%%%%%%%%%%%%%%%%%%%%

% فرمول‌های نمونه:
%$f(x) = \Theta(x)$
%$T(x, y) \leq \log(x) + 2^y + \binom{2n}{n}$
% {\tt code\_function}


% لیست موارد بدون شماره‌گذاری:
%\شروع{فقرات}
%    \فقره مورد اول 
%        \شروع{فقرات}
%            \فقره مورد اول داخلی 
%                \شروع{فقرات}
%                    \فقره اولین مورد خیلی داخلی 
%                    \فقره دومین مورد خیلی داخلی 
%                \پایان{فقرات}
%            \فقره مورد دوم داخلی 
%        \پایان{فقرات}
%    \فقره مورد دوم 
%\پایان{فقرات}

% لیست موارد شماره‌گذاری شده:
%\شروع{شمارش}
%    \فقره یک مورد 
%    \فقره یک مورد دیگر 
%    \فقره باز هم یک مورد دیگر
%\پایان{شمارش}

% برای نوشتن یک فرمول چند خطی که یک شماره بخوره. برای حالت‌های دیگر به سایت زیر مراجعه کنید:
% https://www.overleaf.com/learn/latex/Aligning_equations_with_amsmath
%\begin{align} 
%    \begin{split}
%        (x+y)^3 &= (x+y)^2(x+y) \\
%                &= (x^2+2xy+y^2)(x+y) \\
%                &= (x^3+2x^2y+xy^2) + (x^2y+2xy^2+y^3) \\
%                &= x^3+3x^2y+3xy^2+y^3
%    \end{split}                                 
%\end{align}

% نوشتن ماتریس و آرایه دوبعدی
%\begin{align}
%    A = 
%    \begin{bmatrix}
%        A_{11} & A_{21} \\
%        A_{21} & A_{22}
%    \end{bmatrix}
%\end{align}

% اضافه کردن تصویر با شرح. برای اضافه کردن تصویر بدون شرح، فقط دستور درج‌تصویر را استفاده کنید.
%\شروع{شکل}[h!]
%\تنظیم‌ازوسط
%  \درج‌تصویر[پهنا=0.5\پهنای‌متن]{img.jpg}
%\شرح{شرح تصویر}
%\برچسب{تصویر:۱}
%\پایان{شکل}
%\پایان{زیرمسائل}

%کد زیر را ببینید.
%%\lstinputlisting{source_filename.py}
%\begin{latin}
%\begin{lstlisting}
%# Import the modules
%import sys
%import random
%
%ans = True
%
%while ans:
%question = raw_input("Ask the magic 8 ball a question: (press enter to quit) ")
%
%answers = random.randint(1,8)
%
%if question == "":
%sys.exit()
%
%elif answers == 1:
%print "It is certain"
%
%elif answers == 2:
%print "Outlook good"
%
%elif answers == 3:
%print "You may rely on it"
%
%elif answers == 4:
%print "Ask again later"
%
%elif answers == 5:
%print "Concentrate and ask again"
%
%elif answers == 6:
%print "Reply hazy, try again"
%
%elif answers == 7:
%print "My reply is no"
%
%elif answers == 8:
%print "My sources say no"
%\end{lstlisting}
%\end{latin}
